\section{Discussion}
\label{sec:discussion}

We interpret our findings, explain unexpected results, acknowledge limitations, and discuss broader implications.

\subsection{Key Findings Interpretation}

\subsubsection{Strong Correlation: SA Rules Encode Transferable Knowledge}

Pylint achieves $r=0.87$ correlation with pass@1---far stronger than the $r\geq0.35$ threshold we set for ``moderate'' predictive power. This suggests static analysis rules, originally designed for human-written code, encode expert knowledge that transfers directly to LLM-generated code.

Why does this work? Pylint's $\sim$400 rules capture decades of accumulated wisdom about defect-prone patterns: inconsistent naming, overly complex structures, unused variables, missing documentation. These same anti-patterns appear in incorrect LLM code. The LLM may generate syntactically valid code that violates style conventions \emph{because} it's confused about the underlying logic---style and correctness co-occur.

\textbf{Implication:} SA-based filtering is a cheap, effective quality gate. Running pylint costs $\sim$0.01$\times$ the compute of test execution, yet provides $r=0.87$ predictive power.

\subsubsection{Ensemble Degradation: Redundant Signals}

The weighted ensemble (pylint+radon) achieves $r=0.86$, \emph{below} pylint alone ($r=0.87$). Weight sensitivity analysis shows correlation monotonically increases with pylint weight.

\textbf{Our interpretation:} Both metrics capture overlapping quality dimensions. Pylint rules include complexity checks; radon's cyclomatic complexity measure adds no orthogonal signal. The 0.9/0.1 optimal weighting confirms pylint dominance---any radon contribution introduces noise.

This ``less is more'' finding simplifies deployment: practitioners need track only one metric.

\subsubsection{Cross-Model Variance: GPT-4 Outlier}

All four LLMs show significant SA-correctness correlation (all $p<0.001$), but GPT-4 shows notably lower correlation ($r=0.42$) compared to others ($r\approx0.85$). This outlier inflates variance to $\text{std}=0.19$, exceeding our 0.15 threshold.

\textbf{Possible explanations:}

\begin{enumerate}
\item \textbf{Synthetic data artifact:} Our cross-model analysis used simulated completions. The GPT-4 profile in synthetic data may not reflect real API behavior.

\item \textbf{Qualitatively different code:} GPT-4 may produce code that violates pylint rules but remains functionally correct.

\item \textbf{Model-specific patterns:} Different LLM training data may induce different code style distributions.
\end{enumerate}

\textbf{Most likely interpretation:} Synthetic data methodology artifacts. Real API outputs needed to validate.

\subsection{Limitations}

\textbf{Short Functions Only:} Results based on HumanEval/MBPP problems (30--100 LOC). Repository-level code may show different SA-correctness relationships.

\textbf{Synthetic Multi-Model Data:} Cross-model variance analysis used simulated completions, not real API outputs. The GPT-4 outlier may reflect synthetic data methodology.

\textbf{Mypy Integration Failure:} Mypy errors produced numerical artifacts due to rank-deficient covariance matrix. Type-checking signal remains unexplored.

\textbf{Canonical Solution Bias:} H-M1 correlation computed on canonical solutions where most samples pass. Real LLM outputs have more diverse pass/fail distribution.

\subsection{Broader Impact}

\textbf{Positive Applications:}
\begin{itemize}
\item Cheaper LLM code quality filtering
\item Rapid prototyping feedback
\item Training data curation
\end{itemize}

\textbf{Potential Concerns:}
\begin{itemize}
\item Over-reliance on style metrics
\item Gaming if SA scores become deployment gates
\item False confidence leading to skipped tests
\end{itemize}

\textbf{Mitigation:} SA-based filtering should complement, not replace, test execution.
